\documentclass[11pt]{article}
\usepackage[top=0.75in,bottom=0.85in,left=0.85in,right=0.85in]{geometry}
\usepackage{microtype}
\usepackage{parskip}
\setlength{\parskip}{4pt plus 1pt minus 1pt}
\usepackage{titlesec}
\usepackage{enumitem}   % lists
\usepackage{hyperref}
\usepackage{fancyhdr}
\usepackage{booktabs}
\usepackage{xcolor}
\usepackage{caption}

% --- Compact section spacing ---
\titlespacing*{\section}{0pt}{6pt plus 1pt minus 1pt}{3pt}
\titlespacing*{\subsection}{0pt}{4pt plus 1pt minus 1pt}{2pt}
\titlespacing*{\subsubsection}{0pt}{3pt plus 1pt minus 1pt}{1.5pt}

% --- Lists: deeper indentation ("tabbed over") ---
% Level 1 bullets/numbered:
\setlist[itemize,1]{nosep,leftmargin=2.2em}
\setlist[enumerate,1]{nosep,leftmargin=2.2em}
% Level 2 bullets (nested):
\setlist[itemize,2]{nosep,leftmargin=3.0em}
\setlist[enumerate,2]{nosep,leftmargin=3.0em}

\linespread{1.00}
\captionsetup{font=small,labelfont=bf}

% --- Header/footer (static header avoids title macro issues) ---
\pagestyle{fancy}
\fancyhf{}
\lhead{SITE 6457: TEAM PROPOSAL}
\rhead{Team Proposal}
\cfoot{\thepage}
\setlength{\headheight}{13.6pt}
\setlength{\headsep}{8pt}

% --- Title block (single author block; no \and) ---
\title{Site 6457: A Prisoner's Dilemma -- Team Proposal}
\author{Truman Hess \and Josh Smith \and Larry Wilson \and Wafer Sung \and Michael Markoe}
\date{September 22, 2025}

% --- Avoid widows/orphans ---
\clubpenalty=10000
\widowpenalty=10000

\begin{document}
\maketitle
\vspace{-0.8em}

\section{Basic Gameplay \& Core Mechanics}
\textbf{Core Gameplay Loop.} The player toggles between four prisoners to pathfind and escape a procedurally generated horror environment. Each prisoner must evade the monster, locate the exit, and \textit{stay alive}. The player alternates control to coordinate escapes while non-selected characters are controlled by simple AI. The monster evolves within a round by absorbing traits from captured prisoners, and the next round's difficulty is influenced by how many (and which types of) prisoners were saved previously.

\medskip
\textbf{Unique Gameplay Elements.}
\begin{itemize}
  \item \textit{Monster Evolution}: When the monster captures a prisoner, it gains the strengths and weaknesses of that prisoner. Potential evolutions include speed, pathfinding effectiveness, and environmental manipulation (e.g., doors and entry points).
  \item \textit{Procedural Levels}: Layouts vary each run.
  \item \textit{Unlockable Control}: Potential variant---start as a single prisoner; others become controllable only after they are found.
\end{itemize}

\section{Formal Elements of the Game}

\subsection*{a) Players \& Interaction Modes}
\begin{itemize}
  \item \textbf{Number of players}: MVP single-player; future potential for multiplayer.
  \item \textbf{Interaction style}: MVP toggling between prisoners; future co-op possibilities.
\end{itemize}

\subsection*{b) Objectives}
\begin{itemize}
  \item \textbf{Primary}: Escape with at least one prisoner.
  \item \textbf{Secondary / stretch}: Escape with multiple prisoners.
  \item \textbf{Perfect win}: Escape with all prisoners.
\end{itemize}

\subsection*{c) Procedures}
\begin{itemize}
  \item \textbf{Starting action}: Player spawns alone in a room with a note.
  \item \textbf{Progression}: Explore the generated environment and find other prisoners who can aid in escape. Audio cues indicate monster proximity.
  \item \textbf{Special actions}:
    \begin{itemize}
      \item Each prisoner has a unique skill: (1) speed, (2) environmental manipulation (doors), (3) enhanced hearing.
      \item If the monster catches a prisoner off-screen, the player is notified via sound effects.
      \item On capture, the monster ``trains'' for 5--15 seconds, absorbing that prisoner's skill (e.g., becomes faster after capturing the speed-focused prisoner).
    \end{itemize}
  \item \textbf{Resolving action}: One to four prisoners reach the exit before all are captured.
\end{itemize}

\subsection*{d) Rules}
\begin{itemize}
  \item \textbf{Objects \& concepts}: Characters, monster, environment; the player can move, jump, toggle between characters, and interact with objects (e.g., doors).
  \item \textbf{Restrictions}: The player can command characters to wait or to follow the actively controlled character.
  \item \textbf{Effects}: Non-controlled prisoners generally stay put; depending on attributes, they may hide effectively from the monster.
\end{itemize}

\subsection*{e) Resources}
\begin{itemize}
  \item \textbf{From the monster's perspective}: Prisoners are resources that power the monster up (or down) by absorbing their skills.
  \item \textbf{Exit points}: Initially one exit (may scale with difficulty later).
  \item \textbf{From the player's perspective}: Prisoners are resources---adding prisoners improves team capabilities and score.
\end{itemize}

\subsection*{f) Conflict}
\begin{itemize}
  \item \textbf{Opponent}: The monster.
  \item \textbf{Obstacles}: Doors and room access points.
  \item \textbf{Player dilemma}: If the exit is found before all prisoners are located, decide whether to escape now or keep searching; optionally ``sacrifice'' a prisoner by ordering them to stay hidden.
\end{itemize}

\subsection*{g) Boundaries}
\begin{itemize}
  \item \textbf{Map}: Procedurally generated horror setting (e.g., prison or hospital).
  \item \textbf{Access points}: Some walls are breakable only by certain prisoner types; some doors can be unlocked only by others.
\end{itemize}

\subsection*{h) Outcome}
\begin{itemize}
  \item \textbf{Win}: Highest win is all four escape.
  \item \textbf{Partial win}: One to three escape.
  \item \textbf{Loss}: All are captured by the monster.
\end{itemize}

\section{Narrative / Story Components}
\begin{itemize}
  \item \textbf{Theme / Setting}: Horror atmosphere (hospital or prison).
  \item \textbf{Backstory}: Prisoners are test subjects in an AI training experiment; the monster learns by capturing escapees.
  \item \textbf{Atmosphere Goals}: Excitement and fear via sound and environment; constant trade-offs between personal survival and group success.
\end{itemize}

\section{Work Plan}
Team divides the game into key functional areas and generates tickets on a Trello board.

\subsection*{Functional Breakout}
\begin{itemize}
  \item AI Systems
  \item Procedural Generation
  \item Environmental
  \item Gameplay Control
  \item Audio and UI
\end{itemize}

\subsection*{Key Algorithms / Libraries}
\begin{itemize}
  \item Kruskal's algorithm (minimum spanning tree) to connect rooms.
  \item A* (A-star) for pathfinding.
  \item Behavior Trees for agent control.
\end{itemize}

\end{document}
